% GENERATED FILE. Edit canonical lessons, then run npm run generate:books.
% canonical-source-hash: fnv1a64:2f03b568700baf0f
% canonical-lessons: MR-C384-paristhiti, MR-C384-sthiti, MR-C384-avastha, MR-C384-pec, MR-C384-kasoti

\chapter{Situations}
\label{ch:mr-situations}
\hlchaptermodality{384}

\begin{chapteropening}
\textbf{By the end of this chapter:} \emph{I can describe a situation, a state of things and a dilemma.}

Complete the last lesson of chapter 384: I can describe a situation, a state of things and a dilemma.
\end{chapteropening}

\section[\texorpdfstring{\mr{परिस्थिती} (paristhitī)}{परिस्थिती (paristhitī)}]{\mr{परिस्थिती} (paristhitī) — circumstances; a situation}

\label{lesson:MR-C384-paristhiti}



\begin{quote}
Before the new one: say the Marathi for a trend, then the Marathi for to size up.
\end{quote}

\subsection*{You'll want to know: \mr{परिस्थिती}}
\textbf{\mr{परिस्थिती}} — \emph{paristhitī} — \textquotedblleft{}circumstances; a situation\textquotedblright{}.

\begin{quote}
\mr{घरची} \mr{परिस्थिती} \mr{चांगली} \mr{नव्हती}, \mr{तरीही} \mr{तो} \mr{शिकला}.
\end{quote}

\begin{quote}
\emph{gharcī paristhitī cāṅglī navhatī, tarīhī to śiklā.}
\end{quote}

\begin{quote}
Things were hard at home; even so, he studied.
\end{quote}

Feminine.

Feminine and singular: \mr{घरची} \mr{परिस्थिती}, the situation at home.

\begin{cousinweb}
\textbf{\mr{परिस्थिती}} (\emph{paristhitī}) is a learned word built from Sanskrit \emph{pari-}, \textquotedblleft{}around\textquotedblright{}, and \emph{sthiti}, \textquotedblleft{}standing, state\textquotedblright{} (Monier-Williams): how things stand all round.
\end{cousinweb}

\begin{grammarlens}[title={What you've built}]
One more word for saying what happens if.
\end{grammarlens}

\subsection*{Your turn}
\begin{itemize}
\raggedright
  \item \emph{Say it:} \emph{paristhitī}
  \item \emph{Say it:} \emph{paristhitī}, once more
  \item \emph{Say it:} \emph{gharcī paristhitī cāṅglī navhatī, tarīhī to śiklā.} — Things were hard at home; even so, he studied.
  \item \emph{Recall:} say the Marathi for a trend, then the Marathi for to size up, then say \emph{paristhitī} again
\end{itemize}

\begin{tcolorbox}[breakable,colback=teal!4,colframe=teal!35!black,title={Before you move on}]
What does \mr{परिस्थिती} mean? (\textquotedblleft{}Circumstances; a situation\textquotedblright{}.) Say it once more.
\end{tcolorbox}
\section[\texorpdfstring{\mr{स्थिती} (sthitī)}{स्थिती (sthitī)}]{\mr{स्थिती} (sthitī) — a state; a condition}

\label{lesson:MR-C384-sthiti}



\begin{quote}
Before the new one: say the Marathi for to size up, then the Marathi for circumstances.
\end{quote}

\subsection*{You'll want to know: \mr{स्थिती}}
\textbf{\mr{स्थिती}} — \emph{sthitī} — \textquotedblleft{}a state; a condition\textquotedblright{}.

\begin{quote}
\mr{रस्त्यांची} \mr{स्थिती} \mr{पावसाळ्यात} \mr{खराब} \mr{होते}.
\end{quote}

\begin{quote}
\emph{rastyāñcī sthitī pāvsāḷyāt kharāb hote.}
\end{quote}

\begin{quote}
The state of the roads gets bad in the monsoon.
\end{quote}

Feminine.

Feminine: X \mr{ची} \mr{स्थिती}, the state of X.

\begin{cousinweb}
\textbf{\mr{स्थिती}} (\emph{sthitī}) is a tatsama, taken from Sanskrit \emph{sthiti}, \textquotedblleft{}standing upright, continued existence\textquotedblright{} (CDIAL 13769).
\end{cousinweb}

\begin{grammarlens}[title={What you've built}]
One more word for saying what happens if.
\end{grammarlens}

\subsection*{Your turn}
\begin{itemize}
\raggedright
  \item \emph{Say it:} \emph{sthitī}
  \item \emph{Say it:} \emph{sthitī}, once more
  \item \emph{Say it:} \emph{rastyāñcī sthitī pāvsāḷyāt kharāb hote.} — The state of the roads gets bad in the monsoon.
  \item \emph{Recall:} say the Marathi for to size up, then the Marathi for circumstances, then say \emph{sthitī} again
\end{itemize}

\begin{tcolorbox}[breakable,colback=teal!4,colframe=teal!35!black,title={Before you move on}]
What does \mr{स्थिती} mean? (\textquotedblleft{}A state; a condition\textquotedblright{}.) Say it once more.
\end{tcolorbox}
\section[\texorpdfstring{\mr{अवस्था} (avasthā)}{अवस्था (avasthā)}]{\mr{अवस्था} (avasthā) — a plight; a stage}

\label{lesson:MR-C384-avastha}



\begin{quote}
Before the new one: say the Marathi for circumstances, then the Marathi for a state.
\end{quote}

\subsection*{You'll want to know: \mr{अवस्था}}
\textbf{\mr{अवस्था}} — \emph{avasthā} — \textquotedblleft{}a plight; a stage\textquotedblright{}.

\begin{quote}
\mr{पुरानंतर} \mr{गावाची} \mr{अवस्था} \mr{वाईट} \mr{झाली}.
\end{quote}

\begin{quote}
\emph{purānantar gāvācī avasthā vāīṭ jhālī.}
\end{quote}

\begin{quote}
After the flood the village was in a bad way.
\end{quote}

Feminine.

Also a stage of life: \mr{बालावस्था}, childhood.

\begin{cousinweb}
\textbf{\mr{अवस्था}} (\emph{avasthā}) is a tatsama, taken from Sanskrit \emph{avasthā}, \textquotedblleft{}stability, state, condition\textquotedblright{} (CDIAL 872).
\end{cousinweb}

\begin{grammarlens}[title={What you've built}]
One more word for saying what happens if.
\end{grammarlens}

\subsection*{Your turn}
\begin{itemize}
\raggedright
  \item \emph{Say it:} \emph{avasthā}
  \item \emph{Say it:} \emph{avasthā}, once more
  \item \emph{Say it:} \emph{purānantar gāvācī avasthā vāīṭ jhālī.} — After the flood the village was in a bad way.
  \item \emph{Recall:} say the Marathi for circumstances, then the Marathi for a state, then say \emph{avasthā} again
\end{itemize}

\begin{tcolorbox}[breakable,colback=teal!4,colframe=teal!35!black,title={Before you move on}]
What does \mr{अवस्था} mean? (\textquotedblleft{}A plight; a stage\textquotedblright{}.) Say it once more.
\end{tcolorbox}
\section[\texorpdfstring{\mr{पेच} (pec)}{पेच (pec)}]{\mr{पेच} (pec) — a tight spot; a dilemma}

\label{lesson:MR-C384-pec}



\begin{quote}
Before the new one: say the Marathi for a state, then the Marathi for a plight.
\end{quote}

\subsection*{You'll want to know: \mr{पेच}}
\textbf{\mr{पेच}} — \emph{pec} — \textquotedblleft{}a tight spot; a dilemma\textquotedblright{}.

\begin{quote}
\mr{कोणता} \mr{पर्याय} \mr{निवडावा}, \mr{हा} \mr{मोठा} \mr{पेच} \mr{होता}.
\end{quote}

\begin{quote}
\emph{koṇtā paryāy nivḍāvā, hā moṭhā pec hotā.}
\end{quote}

\begin{quote}
Which option to choose was a real dilemma.
\end{quote}

Masculine.

\mr{पेचात} \mr{पडणे}, to land in a fix; also a screw thread, a twist.

\begin{cousinweb}
\textbf{\mr{पेच}} (\emph{pec}) is a loan from Persian \emph{pech}, \textquotedblleft{}a twist, a coil\textquotedblright{}: a problem that turns back on itself.
\end{cousinweb}

\begin{grammarlens}[title={What you've built}]
One more word for saying what happens if.
\end{grammarlens}

\subsection*{Your turn}
\begin{itemize}
\raggedright
  \item \emph{Say it:} \emph{pec}
  \item \emph{Say it:} \emph{pec}, once more
  \item \emph{Say it:} \emph{koṇtā paryāy nivḍāvā, hā moṭhā pec hotā.} — Which option to choose was a real dilemma.
  \item \emph{Recall:} say the Marathi for a state, then the Marathi for a plight, then say \emph{pec} again
\end{itemize}

\begin{tcolorbox}[breakable,colback=teal!4,colframe=teal!35!black,title={Before you move on}]
What does \mr{पेच} mean? (\textquotedblleft{}A tight spot; a dilemma\textquotedblright{}.) Say it once more.
\end{tcolorbox}
\section[\texorpdfstring{\mr{कसोटी} (kasoṭī)}{कसोटी (kasoṭī)}]{\mr{कसोटी} (kasoṭī) — a test; a touchstone}

\label{lesson:MR-C384-kasoti}



\begin{quote}
Before the new one: say the Marathi for a plight, then the Marathi for a tight spot.
\end{quote}

\subsection*{You'll want to know: \mr{कसोटी}}
\textbf{\mr{कसोटी}} — \emph{kasoṭī} — \textquotedblleft{}a test; a touchstone\textquotedblright{}.

\begin{quote}
\mr{कठीण} \mr{काळ} \mr{हीच} \mr{मैत्रीची} \mr{खरी} \mr{कसोटी} \mr{असते}.
\end{quote}

\begin{quote}
\emph{kaṭhīṇ kāḷ hīc maitrīcī kharī kasoṭī aste.}
\end{quote}

\begin{quote}
Hard times are the true test of friendship.
\end{quote}

Feminine.

\mr{कसोटीला} \mr{उतरणे}, to stand the test.

\begin{cousinweb}
\textbf{\mr{कसोटी}} (\emph{kasoṭī}) is a tadbhava: CDIAL lists Marathi \emph{kasoṭī}, \textquotedblleft{}touchstone\textquotedblright{}, under Sanskrit \emph{kaṣapaṭṭikā}, \textquotedblleft{}touchstone\textquotedblright{} (CDIAL 2973).
\end{cousinweb}

\begin{grammarlens}[title={What you've built}]
One more word for saying what happens if.
\end{grammarlens}

\subsection*{Your turn}
\begin{itemize}
\raggedright
  \item \emph{Say it:} \emph{kasoṭī}
  \item \emph{Say it:} \emph{kasoṭī}, once more
  \item \emph{Say it:} \emph{kaṭhīṇ kāḷ hīc maitrīcī kharī kasoṭī aste.} — Hard times are the true test of friendship.
  \item \emph{Recall:} say the Marathi for a plight, then the Marathi for a tight spot, then say \emph{kasoṭī} again
\end{itemize}

\begin{tcolorbox}[breakable,colback=teal!4,colframe=teal!35!black,title={Before you move on}]
What does \mr{कसोटी} mean? (\textquotedblleft{}A test; a touchstone\textquotedblright{}.) Say it once more.
\end{tcolorbox}
